\begin{claim}{C.56}{Proved}{A uniform word bound for diagonal one-parameter affine groups}
For real weights $\alpha,\beta$, use the group
\[
 N_{\alpha,\beta}=\mathbb R\ltimes\mathbb R^2,\qquad
 (\tau,v)(\tau',w)
 =(\tau+\tau',v+\operatorname{diag}(e^{\alpha\tau},
                                     e^{\beta\tau})w).
                                                        \tag{C.56.1}
\]
For every continuous path $k:[0,t]\to N_{\alpha,\beta}$
with $k(0)=1$, put $P=k([0,t])\cup k([0,t])^{-1}$ and
$g=k(t)$. With the explicit uniform integer
\[
                         M=13\,363\,205,                \tag{C.56.2}
\]
every exponential fraction $\exp(r\log g)$, $0\leq r\leq1$,
belongs to $P^M$. The bound is independent of the
weights and the path. Zero, equal and opposite
weights are included.

In the plane-loop setting of C.54--C.55, let $U^\sharp$
be the union of $U^\ddagger$ and every C.54 orbit whose
actual generated source group has dimension at most
three, and put $X^\sharp=\overline{U^\sharp}$.
The specified exponential contractions extend to a
jointly continuous based contraction on $X^\sharp$,
with jointly continuous inverse and exact
$H$-commutation. It extends the prescribed family on
$X^\ddagger$ and retains the exponential value on
every nonopen orbit.
The finite-representative extension of C.55 also
holds relative to $X^\sharp$, for representatives
outside this larger set with compact quotients
by their full local centralizers.
\end{claim}

\begin{proof}
\emph{The anchor and the exact product error.}
Write $k(s)=(\tau_s,v_s)$ and
$A=\max_{0\leq s\leq t}|\tau_s|$.
For $A=0$, the path is in the translation plane,
and C.37's three-letter additive path identity applies.
For $A>0$, choose a path value or its inverse
$h\in P$ whose first coordinate is $A$.
Conjugate all coordinates by a single translation
so that every nonzero-weight coordinate of $h$
has translation part zero. Its scale in that
coordinate differs from one, so this is possible.
Leave any zero-weight translation coordinate
unchanged. This common conjugacy adds no letters
to a word and preserves all group identities.

Put $\alpha_1=\alpha$, $\alpha_2=\beta$,
$x_i=\alpha_i A$, and $D_i=e^{x_i}-1$.
Direct multiplication gives the translation path
\[
 C_s=[h,k(s)]=(0,(D_1(v_s)_1,D_2(v_s)_2)).               \tag{C.56.3}
\]
For a zero weight both contributions to its
commutator coordinate vanish, including the
unchanged translation part of $h$.
Every $C_s$ or inverse costs at most four letters
of $P$. Conjugation by $h^j$ multiplies the
$i$th translation coordinate by $e^{jx_i}$
for every integer $j$.

Apply C.50.3's three-dimensional additive path
identity to $s\mapsto(\tau_s,v_s)$. It supplies
$m\leq5$, signs $\varepsilon_j\in\{1,-1\}$ and
times $s_j\in[0,t]$ with
\[
 r\tau_t=\sum_{j=1}^{m}\varepsilon_j\tau_{s_j},
 \qquad rv_t=\sum_{j=1}^{m}\varepsilon_jv_{s_j}.
                                                        \tag{C.56.4}
\]
Form the actual product $w=\prod_{j=1}^{m}k(s_j)^{\varepsilon_j}$.
For $\varepsilon_j=1$ put
$p_j=\sum_{\ell<j}\varepsilon_\ell\tau_{s_\ell}$;
for $\varepsilon_j=-1$ put
$p_j=\sum_{\ell\leq j}\varepsilon_\ell\tau_{s_\ell}$.
The full inverse formula in (C.56.1) proves
\[
 w=\left(r\tau_t,
       \left(\sum_{j=1}^{m}\varepsilon_j
           e^{\alpha_i p_j}(v_{s_j})_i\right)_{i=1,2}\right),
              \qquad |p_j|\leq5A.                     \tag{C.56.5}
\]
In particular, the additive identity has not
been substituted for a noncommutative product.

Define
\[
 \theta(z,r)=
 \begin{cases}(e^{rz}-1)/(e^z-1),&z\ne0,\\r,&z=0.\end{cases}
                                                        \tag{C.56.6}
\]
C.54.11 gives
$g^r=(r\tau_t,(\theta(\alpha_i\tau_t,r)(v_t)_i)_i)$.
Put $q=\tau_t/A$, $q_j=p_j/A$, so $|q|\leq1$ and
$|q_j|\leq5$. For a real variable $x$, define the
smooth functions
\[
 \begin{split}
 F_0(x)&=\frac{\theta(qx,r)-r}{e^x-1},\\
 F_j(x)&=-\varepsilon_j\frac{e^{q_jx}-1}{e^x-1}
                                      \quad(1\leq j\leq m),
 \end{split}                                            \tag{C.56.7}
\]
with their continuous values at zero.
By (C.56.4)--(C.56.5), the translation taking
$w$ to $g^r$ by left multiplication has
$i$th coordinate exactly
\[
 F_0(x_i)D_i(v_t)_i+
          \sum_{j=1}^{m}F_j(x_i)D_i(v_{s_j})_i.          \tag{C.56.8}
\]
Terms with $D_i=0$ vanish, just as the original
corrections in that coordinate vanish.

\emph{Uniform analytic bounds at every weight.}
Use $\sigma(x)=\int_0^1e^{ux}\,du$ from C.54.9.
For every real $x$,
\[
 \sigma(x)\geq e^{-|x|},\qquad
                       0\leq\sigma'(x)/\sigma(x)\leq1. \tag{C.56.9}
\]
For $|b|\leq5$,
$(e^{bx}-1)/(e^x-1)=b\sigma(bx)/\sigma(x)$
is smooth at zero, has absolute value at most
$5e^{6|x|}$, and has derivative bounded by
$30e^{6|x|}$.
Indeed its logarithmic derivative when nonzero
is the difference of a number of absolute
value at most $|b|$ and one in $[0,1]$.
The zero function for $b=0$ satisfies the
same bounds directly.

For $F_0$, regard $\theta(z,r)$ as the
probability of $[0,r]$ for the density
$e^{zu}/\sigma(z)$ on $[0,1]$.
Differentiating the two integrals gives
\[
 |\partial_z\theta(z,r)|\leq1,\qquad
 |\partial_z^2\theta(z,r)|\leq6.                         \tag{C.56.10}
\]
To check the second bound explicitly, its
expression in expectations for this density is
\[
 \begin{split}
 &\mathbb E(\mathbf1_{[0,r]}u^2)
 -2\mathbb E(\mathbf1_{[0,r]}u)\mathbb E(u)\\
 &\qquad
 -\mathbb E(\mathbf1_{[0,r]})\mathbb E(u^2)
 +2\mathbb E(\mathbf1_{[0,r]})\mathbb E(u)^2 .
 \end{split}
\]
All expectations lie in $[0,1]$, so the
absolute sum is at most six.
The first derivative is a difference of
two numbers in $[0,1]$.
Consequently
\[
 F_0(x)=
 \frac{q\int_0^1\partial_z\theta(uqx,r)\,du}{\sigma(x)}
                                                        \tag{C.56.11}
\]
has absolute value at most $e^{|x|}$.
The derivative of its numerator is at most
$6q^2\int_0^1u\,du\leq3$ in absolute value.
Equation (C.56.9) therefore bounds
$|F_0'(x)|$ by $4e^{|x|}$.
Every function in (C.56.7) thus satisfies
\[
                 |F(x)|\leq5e^{6|x|},\qquad
                 |F'(x)|\leq30e^{6|x|}.                \tag{C.56.12}
\]

\emph{Bounded interpolation at the two diagonal nodes.}
Set
\[
 \begin{split}
 A_8(x)&=(e^x+e^{-x})^8,\\
 B_8(x)&=(e^x-e^{-x})(e^x+e^{-x})^7,\qquad
 u(x)=\tanh x=B_8(x)/A_8(x).
 \end{split}                                            \tag{C.56.13}
\]
For $f(x)=F(x)/A_8(x)$, the estimates just proved give
\[
 |f(x)|\leq5,\quad |f'(x)|\leq70e^{-2|x|},\quad
        u'(x)=\operatorname{sech}^2x\geq e^{-2|x|}.
                                                        \tag{C.56.14}
\]
We used $A_8(x)\geq e^{8|x|}$ and
$A_8'(x)/A_8(x)=8\tanh x$.
Thus $f$ as a function of $u\in(-1,1)$ is
Lipschitz with constant seventy.

For distinct $x_1,x_2$ put
\[
 b=\frac{f(x_1)-f(x_2)}{u(x_1)-u(x_2)},\qquad
                         a=f(x_1)-bu(x_1).
\]
For equal nodes put $b=0$, $a=f(x_1)$.
In either case,
\[
 |a|\leq75,\quad |b|\leq70,\qquad
 F(x_i)=aA_8(x_i)+bB_8(x_i)\quad(i=1,2).                \tag{C.56.15}
\]
This includes a zero node and coincident weights.
The two node values come from the same function
$F$ in (C.56.7), as follows from the single
logarithmic scale parameter $\tau$.
No uncontrolled divided difference is used.

\emph{Realization by two explicit translation paths.}
Expand the Laurent polynomials
\[
 (z+z^{-1})^8=\sum_{j=-8}^{8}a_jz^j,\qquad
 (z-z^{-1})(z+z^{-1})^7=\sum_{j=-8}^{8}b_jz^j.           \tag{C.56.16}
\]
Their coefficients are integers with
$\sum|a_j|=256$ and $\sum|b_j|\leq256$.
Define continuous translation paths
\[
 \mathcal A_s=\prod_{j=-8}^{8}h^jC_s^{a_j}h^{-j},
 \qquad
 \mathcal B_s=\prod_{j=-8}^{8}h^jC_s^{b_j}h^{-j}.         \tag{C.56.17}
\]
They start at the identity.
Their $i$th translation coordinates are
$A_8(x_i)D_i(v_s)_i$ and $B_8(x_i)D_i(v_s)_i$,
respectively. All factors are translations, so
their order of multiplication is immaterial.

A conjugated copy of $C_s^{\pm1}$ costs at most
$2|j|+4$ letters of $P$. Repetition realizes
each integer coefficient. Hence each value or
inverse of either path costs at most
\[
                       (2\cdot8+4)\,256=5120.          \tag{C.56.18}
\]
The planar additive path identity expresses a
fraction in $[0,1]$ of $\mathcal A_s$ by three
signed values of $\mathcal A|_{[0,s]}$, and
likewise for $\mathcal B$.
For $|a|\leq75$ use seventy-five copies of the
fraction $|a|/75$, with inverse for a negative
sign. For $|b|\leq70$ use seventy copies.
Thus the translation
$\mathcal A_s^a\mathcal B_s^b$ realizes (C.56.15)
times the vector $(D_i(v_s)_i)_i$ and costs at most
\[
                       3(75+70)\,5120=2\,227\,200.     \tag{C.56.19}
\]
The word parameters need not be chosen
continuously: the translation paths to which
the additive identity applies are the actual
continuous paths (C.56.17).

Apply this to $F_0$ at $s=t$ and to every
$F_j$ at $s=s_j$. At most six such translations
give exactly (C.56.8). Their product followed
by $w$ is $g^r$, with length at most
\[
                       5+6\cdot2\,227\,200
                                         =13\,363\,205. \tag{C.56.20}
\]
Identity path values pad shorter words.
This proves (C.56.2), including all zero weights.

\emph{Exact algebraic verification.}
The executed program
\path{code/root/research_data/c56_diagonal_word_certificate.py}
expands the Laurent polynomials, verifies the
interpolation identities and coefficient bounds,
and checks the complete product correction with
rational arithmetic for every signed pattern
of length one through five in six weight regimes.
It also verifies the commutator and all Laurent
conjugation sums in those cases.
The regimes include equal, unequal, opposite and
zero weights. The analytic estimates and signed
path-selection result have their proofs above
and in C.50.3; these finite checks are not used
as their proof.

\emph{Accumulation and the relative extension.}
On $U^\ddagger$, C.55 gives the prescribed source
fractions with bound $1913$. For a new C.54 orbit
of dimension at most three, its actual source is
$\mathbb R^2$, $\operatorname{Aff}^+(\mathbb R)$,
or (C.56.1). The bounds three and (C.56.2)
therefore apply in its specified simply connected
source group. C.54 already proves its individual
closure action and every global quotient equation.
Operator conjugation transports these groups and
their based paths. Consequently the fractions
agree wherever definitions overlap.

Use $M$ in the compact global word-isotopy family
C.50.15. It is a finite union of continuous
images of a compact finite cube, for maps and
inverses. The metabelian identity holds on
$U^\sharp$ and hence on $X^\sharp$.
Every nonopen transitive image there is therefore
one of the line translation, line affine, or
circle rotation groups classified in C.37.9.

For every positive integer $n$, a sequence of
source fractions has the exact residual
\[
 \exp\bigl((nr_i-\lfloor nr_i\rfloor)
                         \log\widetilde f_{t_i}\bigr)
\]
in the same specified simply connected source.
Its exponent lies in $[0,1)$, so the residual
has the same uniform bound $M$.
Compact word parameters give bounded target
residual logarithms on a line.
On a circle, the source residual loop is null
in that same source group; its two-point winding
passes to the limit and is zero.
The real equation is exactly C.50.19, with $M$
in its bound. Division by $n$ identifies the
limiting fraction, including the integer part
of the based circle angle. Every word fixes
a fixed target point.
These are all inputs of C.50's dense-set limit
and diagonal arguments. They prove joint
continuity on $X^\sharp$.

Apply the same proof to the inverse loop.
The resulting family is the orbitwise inverse,
so it is jointly continuous and each map is
a homeomorphism of $X^\sharp$.
The based-edge and operator identities extend
from the dense $U^\sharp$.
The fractions agree with the old family on
$U^\ddagger$; continuity proves agreement
on $X^\ddagger$.

Finally, representatives outside $X^\sharp$
are outside $X^\ddagger$, so the intrinsic
construction and direct local-equivariant-map
limit C.55.7--C.55.12 apply with the same
compact fundamental sets and full global
quotient equations.
Their closure meets $X^\sharp$ only in
nonopen orbits. The actual image on
$X^\sharp$ is metabelian, so the hybrid
value on the intersection is exactly its
prescribed based exponential fraction.
Finite closed-set pasting of the maps and
inverses proves the last assertion.

The uniform word proof above uses one
logarithmic scale parameter. It supplies no
such bound for the four-dimensional product
affine case in C.54. Arbitrary remaining
local quotients, infinitely many inequivalent
representatives and the exterior nonopen-region
extension remain parts of C.12.
\end{proof}
